\documentclass[a4paper]{article}
\title{Appendix}
\usepackage[left=2cm,right=2cm,top=3cm,bottom=2cm]{geometry}
\usepackage{fontspec}
\usepackage{xeCJK}
\setCJKmainfont[]{Noto Sans Mono CJK TC}
\setCJKmonofont[]{Noto Sans Mono CJK TC}
\XeTeXlinebreaklocale "zh"
\XeTeXlinebreakskip = 0pt plus 1pt
\usepackage{anyfontsize}
\usepackage{t1enc}
\usepackage{listings}
\usepackage{url}
\lstset{language=C++,basicstyle=\ttfamily}
\renewcommand{\thepage}{Appendix - \arabic{page}}
\begin{document}

\begin{center}
\textbf{\huge 115 學年度資訊學科能力競賽臺南一中校內初選}\\
\vspace{5mm}
\textbf{\huge 115 學年度資訊學科能力競賽臺南一中校內初選\\試題本}\\
\vspace{10mm}
\rule{17cm}{2pt}\\
\vspace{5mm}

\huge 附錄\\
\end{center}

\fontsize{14pt}{20pt}\selectfont

參考公式:
\noindent\fbox{%
\begin{minipage}{\dimexpr\textwidth-2\fboxsep-2\fboxrule\relax}
\raggedright
\begin{itemize}
    \item 和的平方公式：$(a + b)^2 = a^2 + 2ab + b^2$
    \item 差的平方公式：$(a - b)^2 = a^2 - 2ab + b^2$
    \item 平方差公式：$a^2 - b^2 = (a + b)(a - b)$
    \item 若直角三角形兩股長為 $a$、$b$，斜邊長為 $c$，則 $c^2 = a^2 + b^2$
    \item 若圓的半徑為 $r$，圓周率為 $\pi$，則圓面積 $= \pi r^2$，圓周長 $= 2\pi r$
    \item 凸 $n$ 邊形的內角和為 $(n - 2) \times 180^\circ$（其中 $n \ge 3$）
    \item 若一個等差數列的首項為 $a_1$，公差為 $d$，第 $n$ 項為 $a_n$，前 $n$ 項和為 $S_n$，則：\\
          $a_n = a_1 + (n - 1)d$，$\displaystyle S_n = \frac{n(a_1 + a_n)}{2}$
    \item 若一個等比數列的首項為 $a_1$，公比為 $r$，第 $n$ 項為 $a_n$，則 $a_n = a_1 r^{n - 1}$
    \item 一元二次方程式 $ax^2 + bx + c = 0$ 的解為 $\displaystyle x = \frac{-b \pm \sqrt{b^2 - 4ac}}{2a}$
\end{itemize}
\end{minipage}%
}

\noindent 祝各位都能破台，GL \& HF

\end{document}
